\section{Group Memory $K_g$}
To make the coordinator genuinely adaptive, we introduce a slow group-level memory $K_g$ that stores not the raw noise, but the outcome of meaningful crisis episodes. The key idea is the same as for agent-wise experience: weak fluctuations should not be promoted into memory.

Let $E$ index a crisis episode. We define an episode significance gate
\begin{equation}
q_E = \Theta(C_E - C_{\mathrm{crit}}),
\end{equation}
where $C_E$ is an episode-level crisis score and $C_{\mathrm{crit}}$ is a threshold above which the event is considered informative.

The episode score can be taken as a bounded combination of crisis severity, recovery quality, and tail behavior,
\begin{equation}
score_E = q_E \left(
w_1\,\tilde C_E
+w_2\,(1-\tilde R_E)
+w_3\,\tilde T_E
+w_4\,(1-\tilde D_E)
\right),
\end{equation}
where $\tilde C_E$ is normalized crisis severity, $\tilde R_E$ is normalized recovery quality, $\tilde T_E$ is normalized tail span, and $\tilde D_E$ is normalized dispersion. The exact weights are bounded and chosen so that the memory only reacts to structurally meaningful episodes.

The group memory then evolves as a slow exponential filter,
\begin{equation}
K_g^{(E+1)} = (1-\lambda_g)K_g^{(E)} + \lambda_g\,score_E,
\end{equation}
with $0<\lambda_g\ll 1$.

The memory then feeds back into the coordinator:
\begin{equation}
\phi_{\mathrm{gain}}(t)=\mathrm{clamp}\!\left(
\phi_0\left[1+\alpha_D\mathcal{D}(t)+\alpha_L\overline{L}(t)+\alpha_K K_g(t)\right]
\right),
\end{equation}
\begin{equation}
\kappa(t)=\mathrm{clamp}\!\left(
\kappa_0\left[1+\beta_D\mathcal{D}(t)+\beta_L\overline{L}(t)+\beta_K K_g(t)\right]
\right).
\end{equation}
In this form, $K_g$ does not create hard control. It only shifts the soft control law according to what the system has already learned from meaningful episodes.

\selectlanguage{russian}
\section{Групповая память $K_g$}
Чтобы координатор действительно стал адаптивным, мы вводим медленную групповую память $K_g$, которая хранит не шум, а результат значимых кризисных эпизодов. Идея та же, что и для опыта отдельных агентов: слабые флуктуации не должны превращаться в память.

Пусть $E$ — номер кризисного эпизода. Вводим gate значимости
\begin{equation}
q_E = \Theta(C_E - C_{\mathrm{crit}}),
\end{equation}
где $C_E$ — эпизодный кризисный скор, а $C_{\mathrm{crit}}$ — порог, выше которого событие считается информативным.

Сам скор эпизода можно собирать как ограниченную комбинацию силы кризиса, качества восстановления и хвоста задержки,
\begin{equation}
score_E = q_E \left(
w_1\,\tilde C_E
+w_2\,(1-\tilde R_E)
+w_3\,(1-\tilde T_E)
+w_4\,(1-\tilde D_E)
\right),
\end{equation}
где $\tilde C_E$ — нормированная сила кризиса, $\tilde R_E$ — нормированное качество восстановления, $\tilde T_E$ — нормированный хвост, а $\tilde D_E$ — нормированная дисперсия. Веса выбираются так, чтобы память реагировала только на структурно значимые эпизоды.

Далее групповая память эволюционирует как медленный экспоненциальный фильтр,
\begin{equation}
K_g^{(E+1)} = (1-\lambda_g)K_g^{(E)} + \lambda_g\,score_E,
\end{equation}
где $0<\lambda_g\ll 1$.

После этого память возвращается в координатор:
\begin{equation}
\phi_{\mathrm{gain}}(t)=\mathrm{clamp}\!\left(
\phi_0\left[1+\alpha_D\mathcal{D}(t)+\alpha_L\overline{L}(t)+\alpha_K K_g(t)\right]
\right),
\end{equation}
\begin{equation}
\kappa(t)=\mathrm{clamp}\!\left(
\kappa_0\left[1+\beta_D\mathcal{D}(t)+\beta_L\overline{L}(t)+\beta_K K_g(t)\right]
\right).
\end{equation}
В такой форме $K_g$ не создаёт жёсткого управления. Он лишь смещает закон мягкого контроля в сторону того, чему система уже научилась на значимых эпизодах.
\selectlanguage{english}
